\documentclass[11pt]{article}

\usepackage[margin=1in]{geometry}
\usepackage[T1]{fontenc}
\usepackage[utf8]{inputenc}
\usepackage{csquotes}

\usepackage[backend=biber,style=numeric,sorting=none,maxnames=3,minnames=1]{biblatex}
\addbibresource{refs.bib}

\usepackage{hyperref}

\usepackage{bucket}

\hypersetup{
  colorlinks = true,
  linkcolor  = bucketblue,
  citecolor  = bucketblue,
  urlcolor   = bucketblue,
}

\glossentry{halflife}{half-life}%
  {https://github.com/gianyrox/bucket-foundation/tree/dev/papers/retention-cognates/scripts/hlr.py}%
  {The lag $h$ after which the modeled recall probability of a word falls
   to one half; see \Cref{def:hlr}.}
\glossentry{ned}{normalized edit distance}%
  {https://github.com/gianyrox/bucket-foundation/tree/dev/papers/retention-cognates/lean/Bucket/Retention.lean}%
  {Levenshtein distance between two strings divided by the length of the
   longer one; see \Cref{def:ned}.}
\glossentry{cognate}{cognate}%
  {https://github.com/gianyrox/bucket-foundation/tree/dev/papers/retention-cognates/scripts/features.py}%
  {A word whose Wiktionary etymology shares an ancestor form with a
   translation of it in the learner's interface language; see
   \Cref{def:cognate}.}

\title{Cognates and Forgetting: Do Words Close to a Known Language Last Longer in Memory?}
\author{Bucket Foundation}
\date{2026-11-25}

\begin{document}
\maketitle

\begin{abstract}
Words that resemble a translation in a language the learner already knows
might be forgotten more slowly. We test this on 12,854,226 Duolingo
practice events from 115,222 learners in eight language pairs, adding
Wiktionary cognate features (spelling similarity, IPA similarity and a
shared-etymon flag) to half-life regression. Hypotheses, tests and
equivalence bounds were preregistered. Pooled, a shared etymon lengthens
the modeled half-life by a factor of 1.031 (95\% CI 1.029 to 1.034,
Holm $p < 10^{-100}$) and full spelling overlap by 1.015 (CI 1.010 to
1.019). Both effects sit inside the preregistered equivalence bound of a
7\% change. Held-out MAE moves by $-1.0\times10^{-6}$ (CI $-3.6\times10^{-6}$
to $1.6\times10^{-6}$; TOST against $\pm0.001$, $p<10^{-10}$), with a
minimum detectable effect of $3.7\times10^{-6}$, so the null on prediction
is an absence of effect. Per pair, the etymon effect reaches a factor of
1.069 for Spanish speakers learning English. Cognate status is a real
but small predictor of retention in this data; a scheduler gains no
measurable accuracy from it.
\end{abstract}

\section{Introduction}
\label{sec:intro}

A spaced repetition scheduler predicts when a learner will forget a word
and schedules the next review shortly before. Half-life regression
\autocite{settles2016hlr} models that forgetting from practice history
plus one indicator per word, so a new word starts with no word-level
information. If a word's closeness to the learner's interface language
predicts its \term{halflife}{half-life}, a scheduler can use that closeness
from the first review. Laboratory work found that \term{cognate}{cognates}
are learned faster and forgotten less \autocite{degroot2000hard}. We test
the same question on the public Duolingo traces
\autocite{settles2016data}: 12.9 million practice events from 115,222
learners across eight language pairs.

\section{Related work}
\label{sec:related}

Settles and Meeder introduced half-life regression and released the
traces we use \autocite{settles2016hlr,settles2016data}. Their model
reached its lowest error with lexeme indicator features, which absorb
any fixed property of a word, cognate status included, and so cannot
say which property matters. FSRS \autocite{ye2022fsrs} models memory
stability per card and is the scheduler Polingual runs; it carries no
word-level prior either. De Groot and Keijzer measured cognate status,
concreteness and frequency in a paired-associate study with Dutch
speakers and found cognates both easier to learn and slower to forget
\autocite{degroot2000hard}.

\section{Preliminaries}
\label{sec:prelim}

\begin{definition}[Half-life model]
\label{def:hlr}
For a practice event with lag $\Delta$ days since the last practice and a
feature vector $x$, the predicted recall is $\hat p = 2^{-\Delta/\hat h}$
with $\hat h = 2^{\theta^\top x}$, the \term{halflife}{half-life}. A
coefficient $\theta_j$ is the change in $\log_2 \hat h$ per unit of
feature $j$.
\end{definition}

\begin{definition}[Normalized edit distance]
\label{def:ned}
For strings $a, b$, $\mathrm{ned}(a,b) = \mathrm{lev}(a,b)/\max(|a|,|b|)$,
with $\mathrm{lev}$ the Levenshtein distance. Similarity is
$1-\mathrm{ned}$.
\end{definition}

\begin{definition}[Cognate features]
\label{def:cognate}
Let $T(w)$ be the Wiktionary translations of a lemma $w$ into the
interface language. Spelling similarity is
$\max_{t\in T(w)} 1-\mathrm{ned}(\mathrm{fold}(w),\mathrm{fold}(t))$, where
$\mathrm{fold}$ lowercases and removes diacritics. IPA similarity is the
same maximum over the first Wiktionary IPA transcription of each word.
The word is a \term{cognate}{cognate} when some $t \in T(w)$ shares an
etymological ancestor form with $w$, Proto-Indo-European roots excluded,
or one word lists the other as its source.
\end{definition}

\section{Method}
\label{sec:method}

\paragraph{Data.} The Duolingo half-life regression traces
\autocite{settles2016data} hold one row per word per session: recall
rate in that session, lag since the word was last seen, and counts of
prior correct and incorrect recalls. The dataset is released under
CC BY-NC 4.0 at
\url{https://doi.org/10.7910/DVN/N8XJME}; it permits nonprofit research
and publication of derived results with attribution. The raw file stays
out of the repository; \texttt{scripts/download.sh} fetches it and checks
its size.

\paragraph{Features.} \texttt{scripts/features.py} parses each lexeme tag
to a lemma and looks it up in a Wiktionary extract (Kaikki, CC BY-SA):
3.5 million translation rows, 5.9 million etymology rows and IPA for
words in the six languages. English-interface pairs look up translations
of the learned word into English; the three English-learning pairs look
up translations of the English word into the interface language. Words
with no translation get zero similarity and a missing indicator
(\texttt{trans\_missing}); words with no IPA get the same treatment
(\texttt{phon\_missing}).

\paragraph{Models.} Both models share the base features of
\autocite{settles2016hlr} without lexeme indicators: a bias,
$\sqrt{1+\text{correct}}$, $\sqrt{1+\text{incorrect}}$, plus
$\log_2(1+\text{lemma length})$ as a control, since cognates in this data
run longer. The cognate model adds spelling similarity, IPA similarity,
the shared-etymon indicator and the two missing indicators. The loss is
the one in \autocite{settles2016hlr}: squared recall error plus $0.01$
times squared half-life error, half-life clipped to
$[15\,\text{min}, 274\,\text{days}]$, with an L2 penalty of $10^{-4}$,
minimized by L-BFGS. Events are sorted by time; the first 90\% train and
the last 10\% test. AUC uses the rounded session recall as the label.

\paragraph{Hypotheses.} \texttt{PREREG.md}, committed before any
per-pair result was computed, states the hypotheses. Family P, pooled:
H0 that the test MAE difference $D = \mathrm{MAE}_{\text{cog}} -
\mathrm{MAE}_{\text{base}}$ is zero, and H0 that each of the spelling,
IPA and etymon coefficients is zero, each against a two-sided
alternative. Family Q repeats the four hypotheses for each of the eight
pairs. One pooled estimate of $D$ was seen before the file was written,
from a run with a numerical overflow that was discarded; the file
discloses it.

\paragraph{Tests.} Each hypothesis gets a two-sided Wald $z$ test at
$\alpha = 0.05$. Holm's step-down correction runs within family P (four
tests) and within each of the four per-pair families (eight tests each).
Equivalence uses two one-sided tests (TOST) at $\alpha=0.05$ against
bounds of $\pm 0.001$ recall units for $D$ and $\pm 0.1$ $\log_2$
half-life units for a coefficient, a 7\% change in half-life. The
minimum detectable effect at 80\% power is $2.80\,\mathrm{SE}$. A test
that rejects H0 and passes TOST is reported as a small effect; one that
retains H0 and passes TOST as equivalence; one that retains H0 and fails
TOST as inconclusive.

\paragraph{Standard errors.} Coefficient standard errors are
cluster-robust sandwich estimates with learners as clusters:
$\hat V = H^{-1} M H^{-1}$, with $H$ the Hessian of the mean loss by
central differences of the analytic gradient and $M$ the outer product
of per-learner score sums divided by $n^2$. The SE of $D$ is the
learner-clustered SE of a mean over test events, holding the fitted
parameters fixed; with 11.6 million training events, parameter noise
adds little to it. \texttt{PREREG.md} named 200 learner-clustered
bootstrap refits; that run exceeded the memory of the analysis machine,
and the sandwich estimator replaces it. Both target the same
learner-clustered sampling variance.

\section{Results}
\label{sec:results}

\begin{table}[htbp]
  \centering
  \footnotesize
  \setlength{\tabcolsep}{3pt}
  \caption{Held-out fit per language pair, interface language first.
    $\Delta$MAE is cognate minus base, $\times 10^{3}$, with its 95\%
    interval.}
  \label{tab:fit}
  \begin{tabular}{@{}lrrrrrrrl@{}}
    \toprule
    Pair & Events & Learners & Lexemes & MAE base & MAE cog. & AUC base & AUC cog. & $\Delta$MAE \\
    \midrule
    \csname @@input\endcsname results/table_fit.tex
    \bottomrule
  \end{tabular}
\end{table}

\begin{table}[htbp]
  \centering
  \small
  \caption{Pooled coefficients in $\log_2$ half-life, 95\% intervals.}
  \label{tab:coef}
  \begin{tabular}{@{}lll@{}}
    \toprule
    Feature & Base model & Cognate model \\
    \midrule
    \csname @@input\endcsname results/table_coef.tex
    \bottomrule
  \end{tabular}
\end{table}

\begin{figure}[htbp]
  \centering
  \includegraphics[width=0.9\linewidth]{figures/fig_coef.png}
  \caption{Cognate coefficients per language pair with 95\% intervals,
    from \texttt{figures/fig\_results.py}.}
  \label{fig:coef}
\end{figure}

\begin{figure}[htbp]
  \centering
  \includegraphics[width=0.9\linewidth]{figures/fig_mae.png}
  \caption{Change in held-out MAE from adding cognate features. Below
    zero means the cognate model predicts recall better.}
  \label{fig:mae}
\end{figure}

\begin{table}[htbp]
  \centering
  \scriptsize
  \caption{Preregistered tests. Estimates are $\log_2$ half-life
    coefficients, and $\Delta$MAE $\times10^{3}$. $p_{\text{Holm}}$ is
    corrected within its family; MDE at 80\% power; TOST $p$ against the
    preregistered bound.}
  \label{tab:tests}
  \begin{tabular}{@{}llrlrrrrl@{}}
    \toprule
    Pair & Quantity & Est. & 95\% CI & $z$ & $p_{\text{Holm}}$ & MDE & TOST $p$ & Verdict \\
    \midrule
    \csname @@input\endcsname results/table_tests.tex
    \bottomrule
  \end{tabular}
\end{table}

\paragraph{Prediction.} Adding cognate features leaves held-out error
unchanged in every pair (\Cref{tab:fit}, \Cref{fig:mae}). Pooled,
$D = -1.0\times10^{-6}$ ($z=-0.76$, Holm $p=0.45$), and the 90\%
interval lies far inside $\pm0.001$, so TOST establishes equivalence.
The MDE of $3.7\times10^{-6}$ is under a three-hundredth of the bound:
the test had power to see any effect a scheduler would notice. AUC moves
by at most $2\times10^{-4}$ in any pair.

\paragraph{Coefficients.} Pooled, all three cognate coefficients differ
from zero after Holm correction (\Cref{tab:coef}, \Cref{tab:tests}). A
shared etymon multiplies the modeled half-life by $2^{0.0445}=1.031$
(CI 1.029 to 1.034); spelling similarity of one against zero by 1.015
(CI 1.010 to 1.019). IPA similarity carries a small negative coefficient,
$-0.012$ (CI $-0.021$ to $-0.003$), once spelling is in the model. Every
pooled effect sits inside the $\pm0.1$ bound by TOST. Practice history weighs more: a first incorrect recall moves
$\log_2 \hat h$ by $-0.161(\sqrt2-1) = -0.067$.

\paragraph{Per pair.} \Cref{fig:coef} shows the coefficients by pair.
The etymon effect is positive after correction in five of eight pairs
and largest for Spanish speakers learning English, a factor of
$2^{0.096}=1.069$ (CI 1.061 to 1.077), and for Italian speakers learning
English, 1.060; neither interval is inside the bound. It is negative for
English speakers learning Spanish ($-0.011$) and Portuguese speakers
learning English ($-0.039$). Spelling and IPA coefficients change sign
across pairs, and in English to Italian and Portuguese to English they
pull in opposite directions, a sign that the two measures share variance
and split it differently by pair. No per-pair test is inconclusive: each
either rejects H0 or passes TOST.

\paragraph{Polingual logs.} The local Polingual store
(\texttt{bucket.academy\_progress}) holds 232 language-branch rows from 107 accounts, of which 379 cards
from 15 accounts have two or more reviews and 0 review-log entries exist.
The preregistered threshold was 1,000 learners with per-review history.
That is too small for this analysis, and the store keeps card snapshots
without per-review history, so the study uses the Duolingo traces alone
[empirical: local supabase\_db\_bucket-foundation, 2026-11-25,
\texttt{make polingual}].

\paragraph{Scheduling consequence.} In BucketMath the review interval at
a target retention grows with stability
[bm:BucketMath.Learning.Capacity.lastHold\_mono\_stability] and weekly
upkeep falls as stability rises
[bm:BucketMath.Learning.Capacity.upkeep\_antitone]. A positive cognate
coefficient therefore translates into longer first intervals and lower
upkeep for cognates; \Cref{lem:mono} states the same for the fitted
linear score.

\begin{lemma}[Monotone score]
\label{lem:mono}
If $\theta_c \ge 0$ and $x_c \le y_c$ with other features fixed, then
$\log_2 \hat h(x) \le \log_2 \hat h(y)$.
\end{lemma}

\begin{proof}
$\theta^\top y - \theta^\top x = \theta_c (y_c - x_c) \ge 0$. The Lean
form is \texttt{logHalfLife\_mono} in \Cref{app:lean}.
\end{proof}

\section{Calibration}
\label{sec:calibration}

The base model follows the loss, clipping and features of
\autocite{settles2016hlr} without lexeme indicators. Its pooled test MAE
is 0.127 and AUC 0.537; we did not re-run their released code, so
agreement with their reported figures is not checked here. The
sandwich SEs were checked against the discarded 20-replicate bootstrap
for the pooled $D$: SE $1.3\times10^{-6}$ against an interval half-width
of $2.1\times10^{-6}$ there, consistent at 20 replicates. The
cognate features were checked by hand on 12 random English-interface
lemmas: \textit{instrument} and \textit{operación} score spelling
similarity above 0.88 and a shared etymon; \textit{schlafen} shares the
Proto-Germanic ancestor of \textit{sleep}; \textit{leche} and
\textit{cebola} score near zero. Every reported number is regenerated by
\texttt{make analysis figures} [empirical: Duolingo HLR traces
doi:10.7910/DVN/N8XJME, 2026-11-25, \texttt{make analysis}].

\section{Limitations}
\label{sec:limitations}

The study is observational. Cognates differ from other words in
frequency, part of speech and course position; the length control covers
one of these. The etymology table is incomplete: German \textit{Wasser}
lists only its Middle High German ancestor, so the shared-etymon flag
misses some true cognates and the spelling score carries part of their
signal. IPA strings from different language editions use different
conventions, which compresses IPA similarity toward zero. The base model
omits lexeme indicators; with them, the cognate coefficients would be
unidentified. Without lexeme indicators the fitted intercept puts the typical
half-life near 170 days, so predicted recall sits close to one and any
word-level feature has little room to move the error. Function words dominate event counts, so the pooled
estimates weight them most.

\printbibliography

\appendix

\section{Glossary}
\label{app:glossary}

\printglossaryentry{halflife}
\printglossaryentry{ned}
\printglossaryentry{cognate}

\section{Lean listings}
\label{app:lean}

\texttt{lean/Bucket/Retention.lean} defines Levenshtein distance
(\Cref{def:ned}), the cognate threshold test and the linear log
half-life (\Cref{def:hlr}), and proves \Cref{lem:mono}:

\begin{lstlisting}[style=bucketlean]
def logHalfLife (base c x : Int) : Int := base + c * x

theorem logHalfLife_mono {base c x y : Int} (hc : 0 <= c) (hxy : x <= y) :
    logHalfLife base c x <= logHalfLife base c y :=
  Int.add_le_add_left (Int.mul_le_mul_of_nonneg_left hxy hc) base
\end{lstlisting}

\texttt{lake build} output, quoted verbatim.
% voice-ignore-next 5
\begin{verbatim}
[2/4] Built Bucket.Retention
[3/4] Built Bucket
Build completed successfully (4 jobs).
\end{verbatim}

\end{document}
